\documentclass[10pt,a4paper]{amsart}
\usepackage[margin=19mm]{geometry}
\usepackage{amsmath,amssymb,mathtools}
\usepackage[scr=rsfs,cal=euler]{mathalfa}
\usepackage{xfrac}
\usepackage[unicode,hidelinks,bookmarks=false]{hyperref}
\newcommand{\ie}{i.e.\ }
\newcommand{\cf}{cf.\ }
\newcommand{\resp}{respectively\ }
\newcommand{\Subsection}{\subsection}
\providecommand{\nobreakdash}{\nobreak\hbox{-}}
\setlength{\emergencystretch}{2em}
\multlinegap0pt
\usepackage{mathtools}
\usepackage{amssymb}
\usepackage{bm}
\usepackage{xcolor}
\usepackage{tikz}
\usepackage{float}
\newtheorem{theorem}{Theorem}
\newtheorem{corollary}{Corollary}
\title{Families of Shape-Wilf-Equivalent Claw-Shaped Partially Ordered Patterns}
\author{Sucharita Biswas}
\date{Source: arXiv:2604.27683v1 (CC BY 4.0)}
\begin{document}
\maketitle

\noindent\emph{Source and licence.} Families of Shape-Wilf-Equivalent Claw-Shaped Partially Ordered Patterns. Version arXiv:2604.27683v1 (CC BY 4.0), distributed by arXiv under the Creative Commons Attribution 4.0 International licence, http://creativecommons.org/licenses/by/4.0/, as recorded in the arXiv licence metadata for this exact version and shown on the abstract page https://arxiv.org/abs/2604.27683v1. Full licence notice: \texttt{licenses/arxiv-2604.27683-CC-BY-4.0.txt}.\par\smallskip

\noindent\textit{Editorial scope.} Counted unit: the article's corollary corollary (source0.tex lines 846--849). The article's complete original proof of it is delivered with it (source0.tex lines 851--894, one \texttt{proof} environment of 44 lines). No mathematics is written, repaired or re-derived: the statement and the proof are the article's own lines, in the article's own notation, and no retained line is substituted or re-flowed.  Retained uncounted as the source-local context the proof needs: lines 748--751.  Nothing was omitted from any retained span, so the package carries no omission record.  No retained line contains a citation, so the wrapper carries no bibliography and no bibliography entry.  No retained line contains a figure environment, an image-inclusion command or drawing code, so no image is delivered and the package declares no asset dependency.  Editorial formatting only, in addition: an AMS \texttt{amsart} preamble that restates the article's own declarations and macros used by the retained lines, namely the environments \texttt{theorem}, \texttt{corollary} on separate counters, together with the standard packages those lines need.  The retained environments print in this document as Theorem 1, Corollary 1 in order of appearance, whereas the article prints the same items with the numbering of its own full document.  The complete native source of the article as distributed in the arXiv e-print for this exact version is bundled unchanged as \texttt{sources/199493/source0.tex}.
\begin{theorem}\label{size_I}
    The cardinality of $I$ is 
    $$m - kd - 1 + k \cdot \max(0, d - l + m - 1),$$ which is independent of the boundary parameter $a$ and only depends on $m,k,d,l$.
\end{theorem}

\begin{corollary}
The set $I$ is empty, i.e., $|I|=0$ if and only if 
$a=1$, $m=kd+1$, and $l\ge d(k+1)$.
\end{corollary}

\begin{proof}
From Theorem~\ref{size_I},
$$
|I|=(m-kd-1)+k\cdot \max(0,d-l+m-1),
$$
where both terms are non-negative.

Assume $|I|=0$. Then both terms must be zero. Hence
$
m-kd-1=0 \implies m=kd+1.
$
Using the condition $a+kd\le m$, we get
\[
a+kd\le kd+1 \implies a\le 1.
\]
Since $a\ge 1$, it follows that $a=1$.

Also,
$
\max(0,d-l+m-1)=0
$ implies
$
d-l+m-1\le 0.
$
Substituting $m=kd+1$, we obtain
\[
d-l+kd\le 0
\implies d(k+1)\le l.
\]

Conversely, if $a=1$, $m=kd+1$, and $l\ge d(k+1)$, then
\[
m-kd-1=0
\]
and
\[
d-l+m-1=d(k+1)-l\le 0.
\]
Thus,
\[
k\cdot \max(0,d-l+m-1)=0.
\]
Therefore $|I|=0$.
\end{proof}

\end{document}
